% 第五章 结论与展望
\thesischapterpage{5\quad 结论与展望}

\section{结论与展望}

\subsection{工作总结}
\thesisbodyfont
本文围绕 LEAP Hand 圆柱体连续手内旋转任务展开研究，目标是在低成本灵巧手平台上建立从仿真建模、策略训练到真实部署的完整链路。全文工作可概括为三部分。

在任务定义层面，本文将连续手内旋转落到可评测、可训练的具体场景中。围绕圆柱体对象和竖直轴连续旋转目标，论文明确了系统状态、策略观测、训练期特权信息、相对关节目标动作空间、轴向旋转量计算方式、奖励函数和回合终止条件。这些定义既保留连续旋转任务对旋转效率和抓持稳定性的要求，也使策略输入与真实 LEAP Hand 可读取的本体感觉信息保持一致。

在策略设计层面，本文构建了技能先验驱动的本体自适应旋转框架。第一阶段利用仿真中的特权信息学习旋转技能先验和教师策略，使系统在对象属性已知的训练条件下形成可持续旋转能力；第二阶段冻结动作主干，训练历史编码器从本体感觉历史中恢复同一技能先验，从而得到可部署策略；针对第二阶段在分布外条件下的退化，本文又引入教师——动作一致性约束，并配合更宽的质心偏移训练分布，对部署模块进行定向补强。

在实验验证层面，本文完成了仿真评测和实机部署两部分闭环。仿真中，教师策略能够稳定完成连续旋转，增强后的部署策略在常规分布外评测中取得 \(0.9141\) 的成功率；实机中，同一策略能够在真实 LEAP Hand 上闭环运行，并在标准圆柱体、尺寸变化圆柱体和部分近圆柱真实物体上实现低速受控旋转。与此同时，带棱角饮料瓶的失败也说明，当前方法对“近圆柱”对象仍然更稳，对明显非圆柱对象的泛化还不够。

\subsection{研究结论}
\thesisbodyfont
结合第四章的仿真和实机结果，本文得到以下结论。

第一，圆柱体连续手内旋转可以在低成本 LEAP Hand 平台上通过仿真强化学习获得可用能力。第一阶段教师策略在训练后期能够保持较高的轴向角速度、较低的掉落率和接近完整回合的执行时长，说明当前任务建模和奖励设计能够引导策略学到稳定旋转，而不是只学到短时摆动。

第二，本体历史自适应能够将教师侧能力转化为可部署策略，但对分布外参数变化的耐受度有限。原始第二阶段策略在 ID 条件下已经接近教师表现；当对象质量分布和质心偏移范围向外扩展后，成功率、净旋转圈数和轴向角速度均下降。仅靠本体历史恢复技能先验，尚不足以完全覆盖高风险物理参数带来的偏移。

第三，教师——动作一致性约束是本工作中最有效的增强手段。该约束不改变策略输入形式，而是直接约束部署侧输出贴近教师侧动作语义，因此在 OOD 条件下带来更直接的收益。放宽质心偏移分布也有作用，但它主要补充高风险样本覆盖，并不能单独解决泛化问题。

第四，仿真到真实的差异主要来自执行器、接触条件和观测稳定性，而不是任务目标本身。真实平台的速度上限、通信延迟、关节背隙和接触微滑都会压低实际旋转效率，因此实机结果不应与仿真数值一一对应。即便如此，同一部署策略仍能在真实圆柱体和部分近圆柱物体上形成可重复的闭环旋转，说明本文方法链路具备实际可执行性。

\subsection{后续展望}
\thesisbodyfont
围绕本文已经暴露出的边界，后续工作可从三个方向继续推进。

第一，可将真实执行器特性更细致地纳入仿真。例如，在训练阶段加入电机电流限制、通信延迟、关节背隙和更接近实机的动力学模型，使策略在仿真中提前面对更真实的执行条件，从而降低后续实机迁移落差。

第二，可将实机评测从可执行性判断推进到稳定性度量。目前旋转圈数和角速度主要依赖视觉标记或人工观察，适合做结果判断，但不足以支撑更细的控制分析。后续若加入外部视觉追踪、轻量级姿态标记或多视角估计，实机与仿真的对比会更完整。

第三，可继续扩展对象范围，但这一方向需要更谨慎。带棱角饮料瓶的失败说明，当前策略主要掌握的是圆柱或近圆柱表面的连续接触迁移规律。若要进一步走向多物体泛化，需要同时处理真实动力学建模、对象几何覆盖、姿态测量精度和高风险参数采样等问题，不能仅依赖增加训练轮数。

本文给出了低成本灵巧手实现连续手内旋转的一条可行路径，而更复杂对象上的稳定泛化，仍需在后续工作中继续验证和推进。
